\subsection{积分的性质}

命题 1. --- 对每个正的 \(\mu\) -可积数值函数 \(f\) ，

\[
\int f d | \mu | = \int^ {*} f d | \mu | = \mathrm{N} _ {1} (f) \geqslant 0.\tag{2}
\]

因为，\(\int f d|\mu|\) 与 \(\mathrm{N}_{1}(f)\) 都在 \(L^{1}\) 上连续，且对每个具有紧支集的连续函数 \(f \geqslant 0\) 相等；另一方面，每个函数 \(f \geqslant 0\) （\(L^{1}\) 中的）都是具有紧支集的连续函数 \(\geqslant 0\) 的序列的极限（就平均收敛意义而言）（ \(\S3\) , No.~5, 命题 11）；由此得命题。

推论 1. --- 对每个可积函数 \(\mathbf{f} \in \mathcal{L}_{\mathrm{F}}^{1}\) ，\(|\mathbf{f}|\) 是可积的，且

\[
\int | \mathbf {f} | d | \mu | = \int^ {*} | \mathbf {f} | d | \mu | = \mathrm{N} _ {1} (\mathbf {f}).\tag{3}
\]

我们将经常使用命题 1 及其推论 1，即把 \(\int^{*}fd|\mu|\) 或 \(\mathrm{N}_{1}(f)\) 换成 \(\int fd|\mu|\) ，当所处理的是可积函数 \(\geqslant 0\) 时。例如，为使两个可积函数 \(\mathbf{f},\mathbf{g}\) 等价，必须且只须 \(\int |\mathbf{f} - \mathbf{g}|d|\mu| = 0\) 。

我们回顾，为使函数 \(\mathbf{f}\) 属于 \(\mathcal{L}_{\mathrm{F}}^{p}\) ，必须且只须函数 \(|\mathbf{f}|^{p - 1}\cdot \mathbf{f}\) 属于 \(\mathcal{L}_{\mathrm{F}}^{1}\) （§3, No.~8, 定理 7 的推论 1），即它是可积的；这正是术语« \(p\) 次幂可积函数»的来由。此外：

推论 2. --- 对每个函数 \(\mathbf{f} \in \mathcal{L}_{\mathrm{F}}^{p}\) ，数值函数 \(|\mathbf{f}|^{p}\) 是可积的，且

\[
\mathrm{N} _ {p} (\mathbf {f}) = \left(\int | \mathbf {f} | ^ {p} d | \mu |\right) ^ {1 / p}.\tag{4}
\]

这由 \(|\mathbf{f}|\) 属于 \(\mathcal{L}^p\) （§3, No.~5, 命题 11）这一事实以及公式 (2) 立得。

命题 2. --- 对每个可积函数 f，

\[
\left| \int \mathbf {f} d \mu \right| \leqslant \int | \mathbf {f} | d | \mu |.\tag{5}
\]

由不等式 (1) 通过取极限，并注意到 (3) 以及 \(\mathrm{N}_{1}(\mathbf{f})\) 在 \(L_{F}^{1}\) 上的连续性，立得此式。

定理 1. --- 设 F 与 G 是两个 Banach 空间，u 是 F 到 G 中的一个连续线性映射。对每个取值于 F 的可积函数 f，\(u \circ f\) 是可积的，且

\[
\int \mathbf {u} (\mathbf {f} (x)) d \mu (x) = \mathbf {u} \left(\int \mathbf {f} (x) d \mu (x)\right).\tag{6}
\]

我们已知 \(\mathbf{u} \circ \mathbf{f}\) 是可积的（§3, No.~5, 定理 4）；关系式 (6) 对每个 \(\mathbf{f} \in \mathcal{K}_{\mathrm{F}}\) 成立，故由恒等延拓原理推广到每个可积函数 \(\mathbf{f}\) ：因为 \(\mathbf{f} \mapsto \mathbf{u} \circ \mathbf{f}\) 对于平均收敛拓扑是连续的，这由不等式 \(\mathrm{N}_1(\mathbf{u} \circ \mathbf{f}) \leqslant \| \mathbf{u} \| \cdot \mathrm{N}_1(\mathbf{f})\) 得出。

推论 1. --- 设 \(\mathbf{a}'\) 是 \(\mathbf{F}\) 上的一个连续线性形式。若 \(\mathbf{f}\) 是取值于 \(\mathbf{F}\) 的一个可积函数，则数值函数 \(\langle \mathbf{f}, \mathbf{a}' \rangle\) 是可积的，且

\[
\int \left\langle \mathbf {f} (x), \mathbf {a} ^ {\prime} \right\rangle d \mu (x) = \left\langle \int \mathbf {f} (x) d \mu (x), \mathbf {a} ^ {\prime} \right\rangle .\tag{7}
\]

我们将在 Ch. VI, §1, 习题 7、11 与 12 中看到，可以存在函数 \(\mathbf{f}\) ，取值于一个无穷维 Banach 空间 \(\mathbf{F}\) ，使 \(\langle \mathbf{f}, \mathbf{a}' \rangle\) 对每个连续线性形式 \(\mathbf{a}'\) （\(\mathbf{F}\) 上的）都可积，而 \(\mathbf{f}\) 本身不可积。

推论 2. --- 若 \(\mathbf{a}_k\) （ \(1 \leqslant k \leqslant n\) ）是 F 中的向量，而 \(f_k\) （ \(1 \leqslant k \leqslant n\) ）是可积数值函数，则函数 \(\mathbf{f} = \sum_{k=1}^{n} \mathbf{a}_k f_k\) 是可积的，且

\[
\int \left(\sum_ {k = 1} ^ {n} \mathbf {a} _ {k} f _ {k}\right) d \mu = \sum_ {k = 1} ^ {n} \mathbf {a} _ {k} \int f _ {k} d \mu .\tag{8}
\]

